
\documentclass{article}
\usepackage{colortbl}
\usepackage{makecell}
\usepackage{multirow}
\usepackage{supertabular}

\begin{document}

\newcounter{utterance}

\centering \large Interaction Transcript for game `cladder', experiment `full\_v1.5\_default', episode 420 with qwen3.5{-}4b.
\vspace{24pt}

{ \footnotesize  \setcounter{utterance}{1}
\setlength{\tabcolsep}{0pt}
\begin{supertabular}{c@{$\;$}|p{.15\linewidth}@{}p{.15\linewidth}p{.15\linewidth}p{.15\linewidth}p{.15\linewidth}p{.15\linewidth}}
    \# & \multicolumn{2}{c}{Player} && \multicolumn{2}{c}{Game Master} \\
    \hline

    \theutterance \stepcounter{utterance}  
    & & & \multicolumn{4}{p{0.6\linewidth}}{
        \cellcolor[rgb]{0.9,0.9,0.9}{
            \makecell[{{p{\linewidth}}}]{
                \texttt{\tiny{[P1$\langle$GM]}}
                \texttt{You are an expert in causal inference. The following question is not a typical commonsense query, but rather a meticulously designed question created by a professor specializing in causal inference, intended to assess the students' mastery of the course content.} \\
\texttt{Imagine a self{-}contained, hypothetical world with only the following conditions, and without any unmentioned factors or causal relationships: Vaccination status has a direct effect on vaccination reaction and getting smallpox. Getting smallpox has a direct effect on smallpox survival. Vaccination reaction has a direct effect on smallpox survival. Method 1: We look directly at how vaccination status correlates with smallpox survival in general. Method 2: We look at this correlation case by case according to vaccination reaction. To understand how vaccination status affects smallpox survival, is it more correct to use the Method 1 than Method 2?} \\
\texttt{Start your answer with 'Yes' or 'No', followed by additional reasoning or evidence to support your explanation.} \\
            }
        }
    }
     \\ \\

    \theutterance \stepcounter{utterance}  
    & \multicolumn{4}{p{0.6\linewidth}}{
        \cellcolor[rgb]{0.561,0.737,0.561}{
            \makecell[{{p{\linewidth}}}]{
                \texttt{\tiny{[P1$\rangle$GM]}}
                \texttt{No} \\
\\ 
\texttt{**Reasoning:**} \\
\\ 
\texttt{In causal inference, the goal is to estimate the **causal effect** of a treatment (Vaccination status) on an outcome (Smallpox survival). To do this correctly, we must account for **confounding variables** that influence both the treatment and the outcome.} \\
\\ 
\texttt{Let's analyze the causal structure described:} \\
\texttt{1.  **Vaccination status** \$\textbackslash{}rightarrow\$ **Vaccination reaction**} \\
\texttt{2.  **Vaccination status** \$\textbackslash{}rightarrow\$ **Getting smallpox**} \\
\texttt{3.  **Vaccination reaction** \$\textbackslash{}rightarrow\$ **Smallpox survival**} \\
\texttt{4.  **Getting smallpox** \$\textbackslash{}rightarrow\$ **Smallpox survival**} \\
\\ 
\texttt{Here, **Vaccination reaction** is a **confounder** (or more precisely, a common effect of the treatment that also affects the outcome, creating a specific type of bias known as "collider bias" if we condition on it incorrectly, or simply a mediator if we are trying to estimate the total effect). However, in this specific setup, "Vaccination reaction" is a direct cause of the outcome.} \\
\\ 
\texttt{*   **Method 1 (Total Effect):** This method looks at the raw correlation between Vaccination status and Smallpox survival. This captures the **Total Causal Effect**, which includes both the direct effect of the vaccine on survival (via preventing smallpox) and the indirect effect (via the reaction). This is the standard, correct approach to answer "how vaccination status affects smallpox survival" in its entirety.} \\
\\ 
\texttt{*   **Method 2 (Stratification by Reaction):** This method conditions on "Vaccination reaction." In causal diagrams, conditioning on a variable that is a **descendant** of the treatment (a mediator) blocks the pathway through which the treatment exerts its effect.} \\
\texttt{    *   If we stratify by reaction, we are essentially asking: "What is the effect of vaccination on survival *among those who had a reaction*?" and "What is the effect... *among those who did not*?"} \\
\texttt{    *   By doing this, we sever the link between the vaccination status and the outcome that goes through the reaction. We are no longer measuring the total impact of the vaccine; we are only measuring the effect of the vaccine on survival *conditional* on the reaction having occurred.} \\
\texttt{    *   Furthermore, if "Vaccination reaction" is a confounder (which it is, as it is caused by vaccination and affects survival), conditioning on it in a standard regression without proper adjustment for the causal direction can lead to **M{-}bias** or simply fail to capture the total effect if the goal is the overall population impact. More critically, if the question implies using this to estimate the causal effect, stratifying by a mediator (reaction) prevents us from seeing the full causal chain.} \\
\\ 
\texttt{Therefore, to understand the **overall** effect of vaccination status on survival, Method 1 is the correct approach. Method 2 introduces bias by conditioning on a variable (reaction) that lies on the causal pathway (or is a confounder that, if not handled via a DAG{-}based adjustment strategy rather than simple stratification, distorts the total effect). In the context of "how does X affect Y," the total effect (Method 1) is the primary metric of interest unless the question specifically asks for the effect *conditional* on the reaction.} \\
\\ 
\texttt{Thus, it is **more correct** to use Method 1.} \\
            }
        }
    }
    & & \\ \\

    \theutterance \stepcounter{utterance}  
    & & & \multicolumn{2}{p{0.3\linewidth}}{
        \cellcolor[rgb]{0.9,0.9,0.9}{
            \makecell[{{p{\linewidth}}}]{
                \texttt{\tiny{[GM$|$GM]}}
                \texttt{No} \\
            }
        }
    }
    & & \\ \\

    \theutterance \stepcounter{utterance}  
    & & & \multicolumn{2}{p{0.3\linewidth}}{
        \cellcolor[rgb]{0.9,0.9,0.9}{
            \makecell[{{p{\linewidth}}}]{
                \texttt{\tiny{[GM$|$GM]}}
                \texttt{yes} \\
            }
        }
    }
    & & \\ \\

    \theutterance \stepcounter{utterance}  
    & & & \multicolumn{2}{p{0.3\linewidth}}{
        \cellcolor[rgb]{0.9,0.9,0.9}{
            \makecell[{{p{\linewidth}}}]{
                \texttt{\tiny{[GM$|$GM]}}
                \texttt{game\_result = LOSE} \\
            }
        }
    }
    & & \\ \\

\end{supertabular}
}

\end{document}
